\documentclass[10pt]{ctexart}
\usepackage[a4paper,margin=2cm,top=1.8cm,bottom=1.8cm]{geometry}
\usepackage{amsmath,amssymb,amsthm,bm}
\usepackage[table]{xcolor}
\usepackage{graphicx}
\usepackage{tcolorbox}
\tcbuselibrary{skins,breakable}
\usepackage{enumitem}
\usepackage{array,booktabs}
\usepackage{lastpage}
\usepackage{fancyhdr}
\usepackage{caption}
\usepackage{hyperref}

\definecolor{acc}{HTML}{1B4F8A}
\definecolor{acc2}{HTML}{8A4B1B}
\definecolor{graybg}{HTML}{F2F5F9}
\definecolor{cmpbg}{HTML}{FBF7EC}
\definecolor{cmpline}{HTML}{C9A227}

\xeCJKsetup{CJKmath=true}
\setlist[itemize]{leftmargin=1.4em,itemsep=2pt,topsep=2pt,parsep=0pt}
\setlist[enumerate]{leftmargin=1.9em,itemsep=2pt,topsep=2pt,parsep=0pt}
\setlength{\parindent}{0pt}
\setlength{\parskip}{4pt}
\allowdisplaybreaks
\captionsetup{font=small,labelformat=empty,skip=4pt}

\newcommand{\Rr}{\mathbb{R}}
\newcommand{\eps}{\varepsilon}

\newcommand{\Shead}[2]{%
  \begin{tcolorbox}[enhanced,colback=graybg,colframe=acc,
    boxrule=0pt,leftrule=3pt,arc=2pt,top=3pt,bottom=3pt,left=8pt,right=8pt]
  {\heiti\color{acc}#1}\quad #2
  \end{tcolorbox}}
\newtcolorbox{cmpbox}[1][]{enhanced,breakable,colback=cmpbg,colframe=cmpline,
  boxrule=0pt,leftrule=3pt,arc=2pt,top=2pt,bottom=2pt,left=7pt,right=7pt,#1}

\pagestyle{fancy}\fancyhf{}
\fancyhead[L]{\footnotesize\color{acc}数学分析 · 专题讲义 · 下凸函数}
\fancyhead[R]{\footnotesize\color{acc}配图版}
\fancyfoot[C]{\footnotesize\color{black!55}\thepage\,/\,\pageref{LastPage}}
\renewcommand{\headrulewidth}{0.4pt}

\graphicspath{{figs/out/}}

\begin{document}

\begin{center}
{\LARGE\heiti 下凸函数：从定义到图像}\\[6pt]
{\small 专题讲义（配 6 幅示意图）\quad 对应 2021 期中第 1 题 (3) 与第 4 题 (2)}\\[4pt]
{\footnotesize\color{black!60}先看图，再看定义；每个等价说法都配一张图，最后指出它在试卷里的用法}
\end{center}
\hrule height 1.2pt
\vspace{4pt}

%========================================================
\Shead{一、一句话定义}{弦在曲线上方}

设 $f$ 在区间 $I$ 上有定义。若对任意 $x_1,x_2\in I$ 与任意 $t\in[0,1]$，
\[
f\big(tx_1+(1-t)x_2\big)\ \le\ t f(x_1)+(1-t)f(x_2),
\]
则称 $f$ 在 $I$ 上是\textbf{下凸}（凸）函数。

看图 1：当 $t$ 从 $0$ 变到 $1$，点 $tx_1+(1-t)x_2$ 从 $x_2$ 走到 $x_1$；右端 $t f(x_1)+(1-t)f(x_2)$ 正是对应\textbf{弦}（两点连线）上的高度。不等式说的是——\textbf{弦永远不低于曲线}。

\begin{figure}[h]
\centering
\includegraphics[width=0.92\linewidth]{fig1_chord.png}
\caption{图 1　红色竖段 = 弦高 − 曲线高，它 $\ge 0$ 就是下凸。}
\end{figure}

\textbf{命名的由来}：$y=x^2$ 那样「开口向上」的曲线叫下凸；开口向下（如 $y=-x^2$）的叫\textbf{上凸}（凹）。

%========================================================
\Shead{二、两种凸性的对比}{看弦的位置就分清了}

\begin{figure}[h]
\centering
\includegraphics[width=0.92\linewidth]{fig2_convex_concave.png}
\caption{图 2　同一个「两点连线」操作：左图弦在曲线上方（下凸），右图弦在曲线下方（上凸）。}
\end{figure}

\begin{center}\small
\begin{tabular}{l c c}
\toprule
 & 下凸（凸） & 上凸（凹）\\
\midrule
弦的位置 & 在曲线上方 & 在曲线下方 \\
二阶导（若存在） & $f''\ge 0$ & $f''\le 0$ \\
典型例子 & $x^2,\ e^x,\ |x|$ & $\ln x,\ \sqrt{x},\ -x^2$ \\
\bottomrule
\end{tabular}
\end{center}

\textbf{提醒}：不同教材的「凸 / 凹」叫法可能相反。答题时附上几何描述（弦在曲线上方）或二阶导符号，就不会被误判。

%========================================================
\Shead{三、等价说法三种}{考试任选其一，最好写两种}

下面三种说法与定义\emph{完全等价}，配图见图 4（中点）、图 3（切线）、图 5（二阶导）。

\noindent\textbf{(1) 中点形式（最省事的等价定义）}
\[
f\!\left(\frac{x_1+x_2}{2}\right)\ \le\ \frac{f(x_1)+f(x_2)}{2}.
\]

\begin{figure}[h]
\centering
\includegraphics[width=0.88\linewidth]{fig4_midpoint.png}
\caption{图 4　$x_1=-2,\ x_2=1.6$ 时：曲线在中点的高度 $0.04$，小于弦的中点高度 $3.28$。}
\end{figure}

它只取 $t=\frac12$。由它可以推出 $t=\frac{k}{2^n}$ 的全部情形（反复二分），再对一般 $t$ 取极限，就回到定义。

\noindent\textbf{(2) 切线形式（可导时）}
\[
f(y)\ \ge\ f(x)+f'(x)(y-x)\qquad(\forall x,y\in I).
\]
右边正是曲线在 $x$ 处的切线。所以「下凸」= 曲线始终位于它自己每条切线的\textbf{上方}——图 3 中切线整条压在抛物线下面，只在切点相碰。

\begin{figure}[h]
\centering
\includegraphics[width=0.88\linewidth]{fig3_tangent.png}
\caption{图 3　$f(x)=x^2$ 在 $x_0=0.8$ 处的切线 $y=1.6x-0.64$。}
\end{figure}

\noindent\textbf{(3) 二阶导形式（判别工具）}
若 $f$ 在 $I$ 上二阶可导，则
\[
f\ \text{下凸}\iff f''(x)\ge0\quad(\forall x\in I).
\]
这是考场上最快的一招：算一次二阶导，看符号即可。

\begin{figure}[h]
\centering
\includegraphics[width=0.9\linewidth]{fig5_examples.png}
\caption{图 5　常见函数：抛物线、指数、绝对值都是下凸；对数、根号是上凸。}
\end{figure}

%========================================================
\Shead{四、一个纯几何的推论（常用于证明）}{极值点必是最低点}

下凸函数在内部点处若取极小值，则它就是全局最小值；直观地说，弦都在上方，不可能有「曲线的低洼」在两个点之间还冒出来。这解释了为什么第 2 题 (5) 里「最值不在端点 + Fermat 引理」能直接得出结论。

%========================================================
\Shead{五、它在本次试卷中的用处}{第 4 题 (2) 的常数 $\frac18$ 从哪来}

第 4 题 (2) 要证 $\max|f|\le\frac18(b-a)^2\max|f''|$。证明走到一半会得到
\[
|f(x)|\le\frac{M}{2}(x-a)(b-x),\qquad M=\max|f''| .
\]
接下来的放大就是一道\textbf{上凸}问题：令 $t=\dfrac{x-a}{b-a}\in[0,1]$，则
\[
(x-a)(b-x)=(b-a)^2\,t(1-t)\ \le\ (b-a)^2\cdot\frac14 ,
\]
因为 $t(1-t)$ 是开口向下的抛物线，在 $t=\frac12$（即 $x=\frac{a+b}{2}$）取最大值 $\frac14$（图 6）。

\begin{figure}[h]
\centering
\includegraphics[width=0.9\linewidth]{fig6_oneeighth.png}
\caption{图 6　$t(1-t)$ 的最大值是 $1/4$，乘上前面的 $\frac M2$ 正好得到 $\frac M8$。}
\end{figure}

于是
\[
|f(x)|\le\frac{M}{2}\cdot\frac{(b-a)^2}{4}=\frac18 M(b-a)^2 ,
\]
这就把 $\frac18$ 找出来了。取 $f(x)=\dfrac M2(x-a)(b-x)$ 时等号成立，所以 $\frac18$ 是最优常数。

%========================================================
\Shead{六、小结（背这一页就够）}{}

\begin{cmpbox}
\begin{itemize}
\item \textbf{定义}：$f(tx_1+(1-t)x_2)\le tf(x_1)+(1-t)f(x_2)$，$t\in[0,1]$——弦在曲线上方。
\item \textbf{等价}：中点形式；可导时的切线形式 $f(y)\ge f(x)+f'(x)(y-x)$；二阶导 $f''\ge0$。
\item \textbf{命名}：下凸 $=$ 开口向上 $=$ 弦在上方；上凸 $=$ 凹 $=$ 弦在下方（各书叫法可能相反，答题时补几何描述）。
\item \textbf{考试用法}：判凸性用 $f''$ 符号；证明不等式（如 $\frac18$ 那个常数）时，把二次式化成 $t(1-t)$ 或 $(x-a)(b-x)$ 再用凸性放大。
\end{itemize}
\end{cmpbox}

\end{document}
